\documentclass[sigconf]{aamas}

\usepackage{balance}
\usepackage{amsmath,amsthm}
\usepackage{booktabs}
\usepackage{multirow}
\usepackage{enumitem}
\usepackage{microtype}

\setcopyright{ifaamas}
\acmConference[AAMAS '23]{Proc.\@ of the 22nd International Conference
on Autonomous Agents and Multiagent Systems (AAMAS 2023)}{May 29 -- June 2, 2023}
{London, United Kingdom}{A.~Ricci, W.~Yeoh, N.~Agmon, B.~An (eds.)}
\copyrightyear{2023}
\acmYear{2023}
\acmDOI{}
\acmPrice{}
\acmISBN{}

\acmSubmissionID{CS6002 Course Project}

\title[Sybil-Robust Multi-Item Diffusion Auctions]{Sybil-Robust Multi-Item Diffusion Auctions via Root-Certified Cluster Reduction}

\author{Abhi Jain}
\affiliation{%
  \institution{Indian Institute of Technology Bombay}
  \city{Mumbai}
  \country{India}}
\email{23b0903@iitb.ac.in}

\author{Harshit Raj}
\affiliation{%
  \institution{Indian Institute of Technology Bombay}
  \city{Mumbai}
  \country{India}}
\email{23b0950@iitb.ac.in}

\author{Ridam Jan}
\affiliation{%
  \institution{Indian Institute of Technology Bombay}
  \city{Mumbai}
  \country{India}}
\email{23b1049@iitb.ac.in}

\newtheorem{theorem}{Theorem}
\newtheorem{lemma}[theorem]{Lemma}
\newtheorem{proposition}[theorem]{Proposition}
\newtheorem{corollary}[theorem]{Corollary}
\theoremstyle{definition}
\newtheorem{definition}[theorem]{Definition}
\newtheorem{assumption}[theorem]{Assumption}
\theoremstyle{remark}
\newtheorem{remark}[theorem]{Remark}

\newcommand{\G}{G}
\newcommand{\Ccal}{\mathcal{C}}
\newcommand{\QG}{G/\Ccal}
\newcommand{\rroot}{\rho}
\newcommand{\SW}{\operatorname{SW}}
\newcommand{\Rev}{\operatorname{Rev}}
\newcommand{\util}{\operatorname{u}}
\newcommand{\idom}{\operatorname{idom}}
\newcommand{\argmax}{\operatorname*{arg\,max}}

\begin{abstract}
Diffusion auctions sell items through a social network: buyers not only bid, but also decide whether to forward the auction information. The single-item mechanism of Li et al. resolves the diffusion problem but is vulnerable to fake identities; the Sybil-proof mechanisms of Chen et al. handle false-name manipulation only in the single-item setting; and the multi-item mechanism of Zhao et al. does not address Sybil attacks. We study the open direction at the intersection of these works: Sybil-robust diffusion auctions for multiple homogeneous items. Our contribution is a restricted but rigorous mechanism, \emph{SC-GIDM}, for homogeneous items, single-unit-demand buyers, non-colluding identity-splitting attacks, and a root-certified Sybil-cluster oracle. SC-GIDM contracts each certified cluster to its genuine root and runs GIDM on the quotient graph. We prove that any identity-splitting attack induces only a feasible deviation by the cluster root on the quotient graph, which immediately yields individual rationality, diffusion incentive compatibility, weak budget balance, and Sybil-proofness from the corresponding GIDM guarantees. We also show experimentally that identity-level GIDM can be profitably manipulated, while SC-GIDM removes attacker gain and fake winners under the assumed oracle.
\end{abstract}

% \keywords{diffusion auction, social network, Sybil attack, multi-item mechanism, mechanism design}

\begin{document}
\maketitle

\section{Introduction}
Classical auction models assume that the seller can directly reach all buyers. In a social-network market this assumption fails: the seller initially informs only a local neighborhood, and the auction reaches deeper buyers only if already informed buyers choose to propagate the information. The first diffusion-auction model of Li et al. formalized this setting for a single item and showed that the information diffusion mechanism (IDM) is individually rational (IR), incentive compatible (IC), and weakly budget balanced, while network-VCG may run a deficit \cite{LiHZZ17}. Zhao et al. later generalized this line to multiple homogeneous items via GIDM, again preserving truthful bidding and diffusion incentives under single-unit demand \cite{ZhaoLXHJ18}. However, Chen et al. showed that ordinary diffusion mechanisms can be manipulated by \emph{Sybil attacks}, where one real buyer creates fake identities to alter the graph positions that determine rewards and prices; they proposed the first Sybil-proof single-item mechanisms, STM and SCM, and explicitly left the multi-item case as future work \cite{ChenDWWZ23}. 

This report studies exactly that open direction. Our aim is not to claim a complete solution to arbitrary false-name-proof multi-item mechanism design. Instead, we isolate a setting in which rigorous analysis is possible and the connection to the three baseline papers remains transparent:
\begin{itemize}[leftmargin=1.5em]
    \item the seller holds $K\ge 1$ homogeneous items,
    \item each real buyer has single-unit demand,
    \item Sybil attacks are \emph{non-colluding split attacks}, and
    \item a clustering oracle certifies one genuine root for every Sybil cluster.
\end{itemize}
Within this model we propose \emph{SC-GIDM} (Sybil-Cluster GIDM): contract each certified cluster to its genuine root and run GIDM on the resulting quotient graph.

\subsection{Contributions}
Our report makes four concrete contributions.
\begin{itemize}[leftmargin=1.5em]
    \item We formalize a root-certified cluster model for multi-item diffusion auctions under Sybil attacks.
    \item We design SC-GIDM, a quotient-graph mechanism that combines the cluster idea from SCM with the allocation/payment logic of GIDM.
    \item We prove that, under our assumptions, any identity-splitting attack becomes only a feasible deviation by the cluster root on the quotient graph. Hence SC-GIDM inherits IR, diffusion IC, and weak budget balance from GIDM and is Sybil-proof against such attacks.
    \item We provide a complete worked run on a canonical attacked instance and summarize simulation evidence on random tree and Erd\H{o}s--R\'enyi graphs, showing that attacked GIDM loses real welfare and allocates items to fake identities, while SC-GIDM eliminates attacker gain and fake winners under the oracle model.
\end{itemize}

\subsection{Related Work}
Li et al. introduced the basic model of selling one item over a social network, extended VCG to the network setting, and proposed IDM to restore non-deficit while maintaining truthful diffusion \cite{LiHZZ17}. Zhao et al. studied multiple homogeneous items with unit-demand buyers and proposed GIDM, which preserves IR and IC and guarantees seller revenue at least as good as selling only to the seller's immediate neighbors \cite{ZhaoLXHJ18}. Chen et al. modeled false-name manipulation in diffusion auctions, showed that existing mechanisms such as IDM are Sybil-vulnerable, and proposed STM and SCM for the single-item case \cite{ChenDWWZ23}. Our work sits exactly at the missing intersection of the latter two papers: the multi-item analogue of the Sybil-proof diffusion-auction question.

\section{Model}
Let $\G=(V,E)$ be a reported social graph with seller $s\in V$ and buyers $N=V\setminus\{s\}$. The seller has $K\ge 1$ identical items. Each real buyer $i$ has single-unit demand and value $v_i\ge 0$ for receiving one item. Buyer $i$ reports a bid $b_i$ and a set $r'_i\subseteq r_i$ of neighbors to whom the information is forwarded; thus an action is a pair $(b_i,r'_i)$, as in the baseline diffusion model \cite{LiHZZ17}. A buyer participates only if the seller can reach that buyer through the induced invitation graph.

The usual desiderata are the following.
\begin{definition}[Diffusion incentive compatibility]
A mechanism is diffusion IC if truthful bidding and forwarding to all neighbors is a dominant strategy for every buyer.
\end{definition}
\begin{definition}[Individual rationality]
A mechanism is IR if every truthful buyer obtains nonnegative utility.
\end{definition}
\begin{definition}[Weak budget balance]
A mechanism is weakly budget balanced if $\Rev=\sum_i p_i\ge 0$ for every feasible profile.
\end{definition}

We now extend the model with Sybil attacks.
\begin{definition}[Non-colluding split attack]
A real buyer $i$ may create a finite set $S_i$ of fake identities. The resulting identity set controlled by $i$ is the cluster $C_i=\{i\}\cup S_i$. Fake identities may connect internally and may be attached only through the creator's true neighborhood; distinct real buyers do not collude.
\end{definition}

The buyer's utility is evaluated at the \emph{real-buyer} level. Since demand is unit, buyer $i$ values any allocation that gives at least one item to some identity in $C_i$ by exactly $v_i$.
\begin{equation}
\util_i(x,p)=v_i\,\mathbf{1}[\exists t\in C_i:x_t=1]-\sum_{t\in C_i}p_t.
\label{eq:real-utility}
\end{equation}

The single crucial extra assumption in our theorem-level result is not about incentives but about the defensive layer.
\begin{assumption}[Root-certified Sybil partition]
The mechanism is given a partition $\Ccal=\{C_1,\dots,C_m\}$ of reported identities and a root map $\rroot(C_j)\in C_j$ such that each root is a genuine buyer. Under any split attack by buyer $i$, all fake identities created by $i$ remain in the same cluster as $i$ and do not alter the clusters of other real buyers.
\end{assumption}

We stress that the report does \emph{not} solve Sybil detection. The partition is treated as an external certification layer, mirroring the role of trustworthy non-Sybil structure in SCM \cite{ChenDWWZ23}. This restriction is what allows a mathematically clean reduction.

\section{SC-GIDM}
Given a certified partition $\Ccal$, define the quotient graph $\QG$ as follows: each cluster $C\in\Ccal$ becomes one meta-node; there is an edge $(C,D)$ in $\QG$ if and only if the reported identity graph contains at least one edge from some $u\in C$ to some $v\in D$ with $C\neq D$. 

Our presentation-level sketch used cluster contraction as the main intuition. For theorem-level guarantees we make one further restriction: the meta-bid of a cluster is the bid of its certified root, not the maximum bid among internal identities. This choice makes the macro instance invariant to internal fake bid inflation.

\begin{definition}[SC-GIDM]
For each certified cluster $C$, let its meta-bid be $B_C=b_{\rroot(C)}$. Construct the quotient graph $\QG$ from the reported identity graph. Run GIDM on the meta-instance $(\QG,\{B_C\}_{C\in\Ccal},K)$. If a meta-node $C$ wins an item, allocate the physical item to $\rroot(C)$ and charge the entire meta-payment to $\rroot(C)$.
\end{definition}

The mechanism is deliberately austere: internal fake identities disappear before allocation and pricing, and only the genuine root remains strategically relevant.

\section{Main Results}
We first isolate the key reduction step.

\begin{lemma}[Identity splitting induces only a root deviation]
\label{lem:deviation}
Fix a real buyer $i$ and an arbitrary non-colluding split attack by $i$. After cluster contraction, the attacked identity graph induces a quotient-graph action of cluster $C_i$ that is equivalent to a feasible direct report $(\widehat b_i,\widehat r_i)$ by the certified root $i$, where $\widehat b_i=b'_i$ is the root's own bid and $\widehat r_i\subseteq r_i$.
\end{lemma}

\begin{proof}
Because SC-GIDM uses the certified root bid, internal fake bids are ignored, so the meta-bid is exactly the root's bid $b'_i$. By the non-colluding split-attack model, every external edge of a fake identity in $C_i$ is attached only through a true neighbor of $i$; after contraction, the cluster can therefore reach only a subset of $i$'s true neighbors. Thus the entire split attack becomes a feasible direct deviation by the root on the quotient graph.\qedhere
\end{proof}

The next theorem is the structural core of the paper.
\begin{theorem}[Reduction to GIDM]
\label{thm:reduction}
For a fixed certified partition, SC-GIDM is exactly GIDM run on the quotient graph $\QG$.
\end{theorem}

\begin{proof}
SC-GIDM performs no extra allocation or payment operation beyond: (i) contracting clusters, (ii) replacing each cluster by the certified root bid, and (iii) invoking GIDM on the resulting meta-instance. Hence the output of SC-GIDM is, by definition, the output of GIDM on $\QG$, lifted back to the roots of the winning clusters.\qedhere
\end{proof}

Now we can inherit the known guarantees of GIDM.
\begin{corollary}[IR, diffusion IC, and weak budget balance]
\label{cor:inherit}
Under Assumption~1, SC-GIDM is individually rational, diffusion incentive compatible on the quotient graph, and weakly budget balanced.
\end{corollary}

\begin{proof}
By Theorem~\ref{thm:reduction}, SC-GIDM is just GIDM on a different graph. Zhao et al. proved that GIDM is IR, IC, and weakly budget balanced for homogeneous items and unit-demand buyers \cite{ZhaoLXHJ18}. Applying those guarantees to the quotient graph yields the claim.\qedhere
\end{proof}

The main new property is Sybil-proofness.
\begin{theorem}[Sybil-proofness against non-colluding split attacks]
\label{thm:sybil}
Under Assumption~1, no buyer can increase her real-buyer utility in~\eqref{eq:real-utility} by replacing herself with finitely many fake identities. Equivalently, SC-GIDM is Sybil-proof against non-colluding split attacks.
\end{theorem}

\begin{proof}
Fix a buyer $i$ and any split attack. By Lemma~\ref{lem:deviation}, after contraction the attack induces only a feasible deviation by the root $i$ on the quotient graph. By Corollary~\ref{cor:inherit}, truthful bidding and full diffusion are dominant for that root in the quotient-graph GIDM instance. Therefore the root's utility cannot increase under the induced deviation. Since SC-GIDM allocates and charges only at the root level, the real-buyer utility of $i$ in the attacked identity graph is exactly the quotient-graph utility of the root. Hence the split attack cannot improve real-buyer utility.\qedhere
\end{proof}

\begin{remark}
The theorem does \emph{not} say that arbitrary cluster-based max-bid heuristics are Sybil-proof. The root-bid choice is what makes internal fake bid inflation strategically irrelevant.
\end{remark}

\section{Experimental Framework and Results}
We implemented four mechanisms in code: local-$K$-Vickrey, truthful GIDM, identity-level GIDM under Sybil attack, and SC-GIDM with a perfect cluster oracle. The code uses the same restricted assumptions as the theory: homogeneous items, single-unit demand, and non-colluding split attacks. Two attack generators were implemented; all reported numbers below use the chain attack with the fake bids equal to the attacker's bid. 

The canonical family used in the code has $K=3$ items and one attacker $b_0$ with bid $98$. In the truthful instance, GIDM allocates to $b_1,b_2,b_7$ and buyer $b_0$ gets utility $0$. After one inserted Sybil in a chain attack, identity-level GIDM allocates one item to the fake identity and increases the attacker's utility to $22$, whereas SC-GIDM contracts the fake back into the root cluster and returns the truthful outcome. Full details of this run appear in the appendix.

We also evaluated the mechanisms on 120 random seeds for two graph models: rooted trees and Erd\H{o}s--R\'enyi graphs. Table~\ref{tab:random} reports the averages at $q=0$ and $q=4$ Sybils. Three patterns are stable.
\begin{enumerate}[leftmargin=1.5em]
    \item Attacked GIDM loses \emph{real welfare} as the number of Sybils grows.
    \item Attacked GIDM starts allocating items to fake identities, while SC-GIDM keeps the average number of fake winners at $0$.
    \item Seller revenue under attacked GIDM may even increase, which shows that Sybil robustness is not the same as revenue maximization. SC-GIDM instead preserves the no-attack welfare curve under the oracle model.
\end{enumerate}

\begin{table}[t]
\centering
\caption{Average outcomes over 120 seeds. Welfare is measured over real buyers only.}
\label{tab:random}
\begin{tabular}{llrrrr}
\toprule
Model & Mechanism & $q$ & Welfare & Revenue & Fake winners \\
\midrule
\multirow{4}{*}{Tree} 
& Attacked GIDM & 0 & 345.66 & 289.09 & 0.000 \\
& SC-GIDM       & 0 & 345.66 & 289.09 & 0.000 \\
& Attacked GIDM & 4 & 313.93 & 303.89 & 0.408 \\
& SC-GIDM       & 4 & 345.66 & 289.09 & 0.000 \\
\midrule
\multirow{4}{*}{Erd\H{o}s--R\'enyi} 
& Attacked GIDM & 0 & 351.72 & 305.43 & 0.000 \\
& SC-GIDM       & 0 & 351.68 & 305.43 & 0.000 \\
& Attacked GIDM & 4 & 334.06 & 313.05 & 0.217 \\
& SC-GIDM       & 4 & 351.68 & 305.43 & 0.000 \\
\bottomrule
\end{tabular}
\end{table}

The canonical attacker-gain experiment is even sharper: for $q=1$ and $q=2$ chain Sybils, identity-level GIDM gives attacker gains of $22$ and $13$, while SC-GIDM keeps the gain at $0$ in every tested case. Thus the quotient-graph reduction is not only a proof device; it also captures the mechanism's practical robustness in the simulator.

\section{Conclusions}
The three baseline papers suggest a natural research path: IDM explains truthful diffusion for one item, GIDM extends it to multiple homogeneous items, and STM/SCM show that false identities fundamentally matter. Our report contributes a restricted but rigorous bridge between the latter two directions. Under a root-certified cluster oracle, SC-GIDM contracts each cluster to its genuine root and inherits the IR, diffusion IC, and weak budget balance guarantees of GIDM. The same reduction yields Sybil-proofness against non-colluding split attacks. The experiments support the theorem-level message: identity-level GIDM can be manipulated into fake allocations and lower real welfare, whereas SC-GIDM removes those effects under the assumed oracle. 

The main limitation is also the clearest future-work direction: the clustering layer is assumed rather than derived. Extending the result beyond root-certified partitions, beyond non-colluding split attacks, or beyond homogeneous unit-demand settings remains open.

% \appendix

\section{Full Proofs}

\subsection{Proof of Lemma~\ref{lem:deviation}}
Fix a buyer $i$ and one split attack. Let the attacked identity set be $C_i=\{i\}\cup S_i$. By the definition of SC-GIDM, the cluster bid equals the root bid $b'_i$, so internal fake bids never affect the meta-instance. It remains to show that the cluster's diffusion pattern corresponds to a feasible direct deviation by the root.

Consider any external cluster $D\neq C_i$ adjacent to $C_i$ in the quotient graph. By the construction of the quotient graph, this adjacency is induced by at least one identity-level edge $(u,v)$ with $u\in C_i$ and $v\in D$. Since split attacks may attach fake identities only through the creator's true neighborhood, every external endpoint reachable from some fake identity is in fact a true neighbor of the real buyer $i$. Therefore the set of adjacent clusters reachable from $C_i$ corresponds to a subset of the clusters reachable from $i$'s true neighbor set. This subset is precisely a feasible diffusion report $\widehat r_i\subseteq r_i$. Hence the split attack induces only the feasible direct report $(b'_i,\widehat r_i)$ by the root.

\subsection{Proof of Theorem~\ref{thm:sybil}}
Let $\sigma_i$ be any non-colluding split attack by buyer $i$. By Assumption~1, all identities created by $i$ belong to the same certified cluster and do not alter the certified clusters of other buyers. By Lemma~\ref{lem:deviation}, the identity-level manipulation becomes a feasible direct deviation $(\widehat b_i,\widehat r_i)$ of the root $i$ in the quotient-graph GIDM instance. By Corollary~\ref{cor:inherit}, truthful bidding and full diffusion maximize the root's utility in that quotient-graph instance. Hence
\[
\util_i(\text{truthful on }\QG)\ge \util_i((\widehat b_i,\widehat r_i)\text{ on }\QG).
\]
Because SC-GIDM allocates items and transfers only at the root level, the left-hand side is exactly buyer $i$'s real-buyer utility in the truthful identity graph, and the right-hand side is exactly buyer $i$'s real-buyer utility after executing $\sigma_i$ in the identity graph. Therefore no split attack can improve utility.

\section{A Complete Run on the Canonical Sample}
The code uses the following truthful instance with $K=3$ items. The seller initially informs $b_1$ and $b_3$; the reported edges are
\[
\begin{aligned}
&b_0\!:\{b_2,b_4,b_6\},\quad b_1\!:\{b_3,b_7\},\quad b_2\!:\{b_0,b_7\},\quad b_3\!:\{b_1\},\\
&b_4\!:\{b_0,b_5\},\quad b_5\!:\{b_4\},\quad b_6\!:\{b_0\},\quad b_7\!:\{b_1,b_2\},
\end{aligned}
\]
and the bids are
\[
(b_0,b_1,b_2,b_3,b_4,b_5,b_6,b_7)=(98,34,5,1,19,85,76,61).
\]

\paragraph{Truthful GIDM.}
The implementation outputs the allocation
\[
x(b_1)=x(b_2)=x(b_7)=1,
\]
all other allocations $0$, and payments
\[
p(b_1)=0,\qquad p(b_2)=1,\qquad p(b_7)=0.
\]
Hence the attacker's truthful utility is $0$.

\paragraph{One-chain Sybil attack against identity-level GIDM.}
Buyer $b_0$ creates one fake identity $\mathrm{syb}_{b_0}$ with the same bid $98$, and replaces the outgoing report by the chain
\[
b_0\to \mathrm{syb}_{b_0}\to \{b_2,b_4,b_6\}.
\]
Identity-level GIDM now allocates to $b_1$, $b_7$, and $\mathrm{syb}_{b_0}$, and outputs the nonzero payments
\[
p(b_1)=0,\qquad p(b_7)=0,\qquad p(b_2)=-75,\qquad p(\mathrm{syb}_{b_0})=76.
\]
The real buyer $b_0$ receives one item through the fake identity and pays $76$, so its real-buyer utility becomes
\[
98-76=22.
\]
Thus the split attack is profitable.

\paragraph{SC-GIDM on the attacked instance.}
The oracle returns the certified cluster partition with one nontrivial cluster
\[
C_{b_0}=\{b_0,\mathrm{syb}_{b_0}\},
\]
all other buyers remaining singleton clusters. SC-GIDM replaces this cluster by the certified root $b_0$ with bid $98$. After contraction, the quotient graph is exactly the original truthful graph again. Therefore GIDM on the quotient graph returns the same outcome as the truthful instance:
\[
x(b_1)=x(b_2)=x(b_7)=1,
\]
with the same payments
\[
p(b_1)=0,\qquad p(b_2)=1,\qquad p(b_7)=0.
\]
Hence the attacker's utility returns to $0$. This is the concrete sample run used in the simulation code and summarized in the main text.

\balance
\bibliographystyle{ACM-Reference-Format}
\bibliography{sample}

\end{document}
